\section{Kết luận}
\label{sec:conclusion}

Nhóm đã tự cài đặt và so sánh bốn bộ phân lớp trên MNIST với cùng một giao thức không rò rỉ. MLP cho kết quả tốt nhất (\accTestMLP\% trên test) với chi phí nhỏ; AdaBoost-SAMME gần sát (\accTestADA\%) nhưng đắt hơn nhiều về thời gian huấn luyện và dự đoán; softmax (\accTestSM\%) là lựa chọn nhanh và dễ đọc; Naïve Bayes (\accTestNB\%) là mốc đơn giản nhất và bộc lộ rõ cái giá của giả thiết độc lập. Kết quả của các bản tự cài đặt khớp scikit-learn ở mức cho phép của từng thuật toán, và bản xuất ONNX khớp bản Python.

Bài học thực nghiệm lớn nhất: phần lớn lợi ích đến từ \emph{chọn cấu hình bằng validation một cách có kỷ luật} (ví dụ \texttt{var\_smoothing} cho Naïve Bayes, độ sâu cây cho AdaBoost, hình dạng mạng cho MLP), và từ việc ghi lại trung thực cả những thí nghiệm không rõ ràng (tăng cường dữ liệu chỉ cải thiện nhẹ và chưa đủ để kết luận chắc chắn với ba seed).

\textbf{Hướng tiếp theo.} Chạy nhiều seed và nhiều phép chia cho mọi mô hình để có khoảng tin cậy thật; tìm kiếm tương tác siêu tham số (ví dụ tìm ngẫu nhiên hoặc Bayes) trong cùng ngân sách; thử một mạng tích chập như phần mở rộng tùy chọn; hiệu chỉnh xác suất của AdaBoost; và kiểm tra bản demo trên dữ liệu viết tay thật ngoài MNIST.
